\documentclass[11pt]{article}
\usepackage[margin=1in]{geometry}
\usepackage{amsmath,amssymb,amsthm}
\usepackage{graphicx}
\usepackage[hidelinks]{hyperref}

\newtheorem{theorem}{Theorem}
\newtheorem{lemma}{Lemma}
\newtheorem{definition}{Definition}
\newtheorem{remark}{Remark}

\newcommand{\N}{\mathbb N}
\newcommand{\R}{\mathbb R}
\newcommand{\supp}{\operatorname{supp}}
\newcommand{\dist}{\operatorname{dist}}

\title{A Counterexample to Separable Determination of Fixed-Direction
Sigma-Directional Porosity}
\author{Run \texttt{fa\_banach\_001}, lane 7}
\date{June 16, 2026}

\begin{document}
\maketitle

\section*{Source Question}

The source paper is M. Cuth, M. Rmoutil and M. Zeleny,
\emph{On Separable Determination of Sigma-P-Porous Sets in Banach Spaces},
arXiv:1309.2174.  On page 11, after Remark 4.8, the authors ask:

\begin{quote}
It is an open problem whether the notion of $\sigma$-directional porosity
(see [12] for the definition) is separably determined in the sense of
Corollary 5.11.
\end{quote}

The same passage notes that the definition in Lindenstrauss--Preiss--Tiser
is slightly different and equivalent to the cited definition only in separable
Banach spaces.  The counterexample below addresses the fixed-direction
convention cited from Zajicek's paper [12].

\begin{center}
\includegraphics[width=\textwidth]{figures/open_problem_crop.png}
\end{center}

\section*{Definition Used}

We use the fixed-direction convention indicated by the source paper.  A set
$B$ in a Banach space $X$ is porous in the direction $v\neq 0$ if there is
$0<c<1$ such that, for every $b\in B$ and every $\varepsilon>0$, there is a
scalar $t$ with $0<\|t v\|<\varepsilon$ and
\[
  B_X(b+t v,c\|t v\|)\cap B=\varnothing .
\]
A set is $\sigma$-directionally porous if it is a countable union of sets
which are porous in fixed directions, with the direction allowed to vary from
one member of the countable cover to another.

\section*{Proof Intuition}

The obstruction is that a countable decomposition can use only countably many
directions.  In a nonseparable Hilbert space, countably many direction vectors
miss some coordinate entirely.  On the affine hyperplane where that coordinate
equals one, all proposed directions are tangent.  Therefore each proposed
porous piece is nowhere dense inside that complete affine hyperplane, so
Baire's theorem prevents the pieces from covering it.

Separable subspaces do not see this obstruction.  A separable subspace of
$\ell_2(\Gamma)$ is supported on countably many coordinates, and the example
then becomes a countable union of proper affine hyperplanes, each porous in a
fixed direction inside the subspace.

\begin{theorem}
Let $\Gamma$ be an uncountable set and let
$X=\ell_2(\Gamma)$ over the real scalars.  Define
\[
  A=\{x\in X:\ x_\gamma=1\text{ for some }\gamma\in\Gamma\}.
\]
Then $A$ is closed, hence Suslin; $A$ is not $\sigma$-directionally porous in
$X$; but for every separable linear subspace $V\subset X$, the set
$A\cap V$ is $\sigma$-directionally porous in $V$.
\end{theorem}

\begin{proof}
First, $A$ is closed.  Suppose $x^m\in A$ and $x^m\to x$ in $\ell_2(\Gamma)$.
For each $m$ choose $\gamma_m$ with $x^m_{\gamma_m}=1$.  If
$\|x^m-x\|<1/2$, then $|x_{\gamma_m}|>1/2$.  Since an $\ell_2(\Gamma)$ vector
has only finitely many coordinates of modulus greater than $1/2$, one
coordinate $\gamma$ occurs along an infinite subsequence.  Passing to that
subsequence gives $x_\gamma=1$, so $x\in A$.

We next prove that $A$ is not $\sigma$-directionally porous.  Assume, toward a
contradiction, that
\[
  A=\bigcup_{n=1}^\infty A_n
\]
where each $A_n$ is porous in the fixed direction $v_n\neq 0$.  Each
$v_n\in\ell_2(\Gamma)$ has countable support, hence
\[
  S=\bigcup_{n=1}^\infty \supp v_n
\]
is countable.  Choose $\gamma\in\Gamma\setminus S$ and consider the closed
affine hyperplane
\[
  H_\gamma=\{x\in X:\ x_\gamma=1\}.
\]
Then $H_\gamma\subset A$, and $v_n$ is tangent to $H_\gamma$ for every $n$.

For each $n$, the set $A_n\cap H_\gamma$ is porous in the complete metric
space $H_\gamma$, in the same direction $v_n$.  Indeed, if
$h\in A_n\cap H_\gamma$, then the holes for $A_n$ centered at
$h+t v_n$ have centers still in $H_\gamma$, because $(v_n)_\gamma=0$; their
intersections with $H_\gamma$ are holes for $A_n\cap H_\gamma$.
Such a porous subset is nowhere dense in $H_\gamma$: if some nonempty
relatively open ball in $H_\gamma$ were contained in its closure, taking a
point of the set inside that ball and then a sufficiently small porosity hole
would produce a smaller relatively open ball disjoint from the set, a
contradiction.

Thus $H_\gamma=\bigcup_n(A_n\cap H_\gamma)$ is a countable union of nowhere
dense subsets of the complete metric space $H_\gamma$, contradicting Baire's
theorem.  Therefore $A$ is not $\sigma$-directionally porous in $X$.

It remains to check separable sections.  Let $V\subset X$ be separable.  Choose
a countable dense subset $D\subset V$.  The set
\[
  S_V=\bigcup_{d\in D}\supp d
\]
is countable, and $V\subset \ell_2(S_V)$.  Hence
\[
  A\cap V
  =\bigcup_{\gamma\in S_V}\{x\in V:\ x_\gamma=1\}.
\]
For each $\gamma\in S_V$, the set $\{x\in V:\ x_\gamma=1\}$ is either empty
or a proper closed affine hyperplane of $V$.  If it is nonempty, choose
$w\in V$ with $w_\gamma\neq 0$.  With $\varphi_\gamma(x)=x_\gamma$ on $V$,
any point $x$ in this affine hyperplane satisfies
\[
  \dist_V(x+t w,\{y\in V:\ y_\gamma=1\})
  \ge \frac{|t|\,|\varphi_\gamma(w)|}{\|\varphi_\gamma\|}.
\]
Thus the hyperplane is porous in the fixed direction $w$.  Since $S_V$ is
countable, $A\cap V$ is $\sigma$-directionally porous in $V$.
\end{proof}

\section*{Consequence}

Taking $V_0=\{0\}$, every closed separable subspace $V\supset V_0$ satisfies
that $A\cap V$ is $\sigma$-directionally porous in $V$, while $A$ is not
$\sigma$-directionally porous in $X$.  Hence the Corollary 5.11 style
separable-determination statement fails for the fixed-direction
$\sigma$-directional porosity convention cited in the source question.

\section*{Verification and Scope}

The proof is purely analytic and uses no computation.  The central checks are:
\begin{itemize}
\item $A$ is closed, so it satisfies the source paper's Suslin hypothesis.
\item Countably many proposed directions in $\ell_2(\Gamma)$ miss some
coordinate $\gamma$.
\item Baire's theorem applies to the complete affine Hilbert space
$H_\gamma$.
\item Every separable subspace of $\ell_2(\Gamma)$ is supported on countably
many coordinates.
\end{itemize}

The result is stated for the fixed-direction convention referenced by the
source paper through Zajicek [12].  It is not claimed to settle a different
nonseparable convention in which the direction is allowed to vary point by
point inside a single porous set.

\section*{Novelty Check}

The run indexes were searched for arXiv:1309.2174 and for the phrases
\texttt{sigma-directional porosity}, \texttt{directional porosity}, and
\texttt{separably determined}; no existing packet or proof-gap entry for this
question was found.  The local arXiv source corpus was searched for the same
phrases and found the source question plus related uses of
$\sigma$-directionally porous sets, but no later local answer to this exact
separable-determination problem.  Web searches for exact phrases such as
\texttt{"sigma-directional porosity" "separably determined"} and close variants
did not return an explicit later solution.  Human review should verify the
definition in Zajicek's cited paper [12], because the source paper emphasizes
that conventions differ outside separable spaces.

\begin{thebibliography}{9}
\bibitem{CRZ}
M. Cuth, M. Rmoutil and M. Zeleny,
\emph{On Separable Determination of Sigma-P-Porous Sets in Banach Spaces},
arXiv:1309.2174.

\bibitem{Zajicek2005}
L. Zajicek,
\emph{On $\sigma$-porous sets in abstract spaces},
Abstract and Applied Analysis 2005, no. 5, 509--534.

\bibitem{LPT}
J. Lindenstrauss, D. Preiss and J. Tiser,
\emph{Frechet differentiability of Lipschitz functions and porous sets in
Banach spaces}, Annals of Mathematics Studies 179, Princeton University Press,
2012.
\end{thebibliography}

\end{document}
